% ============================================================================
% CHAPTER 1I: INDUSTRIAL AND COMPUTER SYSTEMS
% ============================================================================

\chapter{Industrial and Computer Systems}
\label{ch:industrial_computer}

This section covers systems from manufacturing, process control, and computer science---domains where control theory meets discrete events, queueing, and information flow.

% ============================================================================
\section{Industrial Process Control}
\label{sec:industrial_process}
% ============================================================================

\subsection{Conveyor Belt System}

\input{figures/fig_ch15_conveyor.tex}

Motor-driven conveyor dynamics:
\begin{equation}
J_{\text{eq}}\dot{\omega} + B\omega = K_t i - \tau_L
\end{equation}

where $J_{\text{eq}}$ includes reflected inertia of belt and load.

Belt velocity: $v = r\omega$ where $r$ is the drive roller radius.

Transfer function (speed control):
\begin{equation}
\frac{V(s)}{I(s)} = \frac{rK_t}{J_{\text{eq}}s + B}
\end{equation}

\subsection{Temperature Control Zone}

Industrial ovens often have multiple zones:

\input{figures/fig_ch15_temp_zones.tex}

Each zone:
\begin{equation}
C_i\frac{\dd T_i}{\dd t} = Q_i - \frac{T_i - T_{i-1}}{R_{i-1,i}} - \frac{T_i - T_{i+1}}{R_{i,i+1}} - \frac{T_i - T_{\infty}}{R_{i,\infty}}
\end{equation}

% ============================================================================
\section{Robotic Manipulators}
\label{sec:robotics}
% ============================================================================

\subsection{Single-Link Robot Arm}

\input{figures/fig_ch15_single_link.tex}

Equation of motion:
\begin{equation}
J\ddot{\theta} + B\dot{\theta} + mg\ell\cos\theta = \tau
\end{equation}

Linearized about horizontal ($\theta_0 = 0$):
\begin{equation}
J\ddot{\theta} + B\dot{\theta} + mg\ell = \tau
\end{equation}

Note: Gravity creates a constant disturbance torque.

\subsection{Two-Link Planar Arm}

The dynamics become coupled and nonlinear:
\begin{equation}
\mM(\vecq)\ddot{\vecq} + \mC(\vecq, \dot{\vecq})\dot{\vecq} + \vecg(\vecq) = \boldsymbol{\tau}
\end{equation}

where:
\begin{itemize}
    \item $\mM(\vecq)$ = mass matrix (configuration-dependent)
    \item $\mC(\vecq, \dot{\vecq})$ = Coriolis and centrifugal terms
    \item $\vecg(\vecq)$ = gravity vector
    \item $\boldsymbol{\tau}$ = joint torques
\end{itemize}

For a two-link arm:
\begin{align}
M_{11} &= m_1 \ell_{c1}^2 + m_2(\ell_1^2 + \ell_{c2}^2 + 2\ell_1\ell_{c2}\cos\theta_2) + I_1 + I_2 \\
M_{12} &= m_2(\ell_{c2}^2 + \ell_1\ell_{c2}\cos\theta_2) + I_2 \\
M_{22} &= m_2\ell_{c2}^2 + I_2
\end{align}

% ============================================================================
\section{CNC Machine Tools}
\label{sec:cnc}
% ============================================================================

\subsection{Axis Drive Model}

Each axis of a CNC machine:
\begin{equation}
J_{\text{axis}}\ddot{x} + B\dot{x} + F_c\sgn(\dot{x}) = F_{\text{motor}} - F_{\text{cut}}
\end{equation}

where $F_c$ is Coulomb friction and $F_{\text{cut}}$ is cutting force disturbance.

\subsection{Servo System}

Position loop transfer function:
\begin{equation}
\frac{X(s)}{X_{\text{cmd}}(s)} = \frac{K_p K_v/s}{1 + K_p K_v/s + s^2/\omega_n^2} \approx \frac{K_p}{\tau_p s + 1}
\end{equation}

for well-tuned systems where $K_v = K_p/\tau_p$.

\subsection{CNC Machine Control Loop}

A complete CNC control architecture uses nested feedback loops for position tracking with real-time disturbance rejection:

\input{figures/fig_ch15_cnc_control_loop.tex}

% ============================================================================
\section{Power Systems}
\label{sec:power_systems}
% ============================================================================

\subsection{Synchronous Generator}

Swing equation:
\begin{equation}
\frac{2H}{\omega_s}\frac{\dd^2\delta}{\dd t^2} + D\frac{\dd\delta}{\dd t} = P_m - P_e
\end{equation}

where:
\begin{itemize}
    \item $H$ = inertia constant (seconds)
    \item $\delta$ = rotor angle (electrical radians)
    \item $\omega_s$ = synchronous speed
    \item $P_m$ = mechanical power
    \item $P_e$ = electrical power
\end{itemize}

Linearized about operating point:
\begin{equation}
\frac{\Delta\delta(s)}{\Delta P_m(s)} = \frac{\omega_s/2H}{s^2 + (D\omega_s/2H)s + K_s\omega_s/2H}
\end{equation}

where $K_s = \partial P_e/\partial\delta$ is the synchronizing coefficient.

\subsection{Automatic Generation Control (AGC)}

Area control error:
\begin{equation}
\text{ACE} = \Delta P_{\text{tie}} + B\Delta f
\end{equation}

where $B$ is the frequency bias setting.

% ============================================================================
\section{Queueing Systems}
\label{sec:queueing}
% ============================================================================

\subsection{M/M/1 Queue}

\input{figures/fig_ch15_mm1_queue.tex}

Steady-state:
\begin{align}
\text{Utilization:} \quad \rho &= \lambda/\mu \\
\text{Mean queue length:} \quad L &= \frac{\rho}{1-\rho} \\
\text{Mean waiting time:} \quad W &= \frac{1}{\mu - \lambda}
\end{align}

\input{figures/fig_ch15_queueing_dynamics.tex}

\subsection{Fluid Flow Model}

For control design, queues can be approximated as fluid:
\begin{equation}
\frac{\dd q}{\dd t} = \lambda(t) - \mu(t)
\end{equation}

where $q$ is queue length, $\lambda$ is arrival rate, and $\mu$ is service rate.

Transfer function:
\begin{equation}
\frac{Q(s)}{\Lambda(s) - M(s)} = \frac{1}{s}
\end{equation}

% ============================================================================
\section{Computer Systems}
\label{sec:computer_systems}
% ============================================================================

\subsection{CPU Load Control}

CPU utilization dynamics:
\begin{equation}
\tau\frac{\dd U}{\dd t} + U = K_u \cdot R
\end{equation}

where $U$ is utilization, $R$ is admission rate, and $\tau$ depends on job service time.

\subsection{Network Congestion Control}

TCP congestion window dynamics (simplified):
\begin{equation}
\frac{\dd W}{\dd t} = \frac{1}{RTT} - \frac{W}{2} \cdot p
\end{equation}

where $W$ is congestion window, $RTT$ is round-trip time, and $p$ is packet loss probability.

\subsection{Memory Management}

Page fault rate model:
\begin{equation}
f = f_0 \left(\frac{W_s}{M}\right)^\alpha
\end{equation}

where $M$ is allocated memory, $W_s$ is working set size, and $\alpha$ is a workload parameter.

% ============================================================================
\section{Data Center Thermal Management}
\label{sec:datacenter}
% ============================================================================

\input{figures/fig_ch15_datacenter_thermal.tex}

Thermal dynamics:
\begin{equation}
C_{\text{room}}\frac{\dd T}{\dd t} = Q_{\text{IT}} - \dot{m}c_p(T - T_{\text{supply}})
\end{equation}

where $\dot{m}$ is air mass flow rate (controlled by CRAC fans).

% ============================================================================
\section{Building Automation}
\label{sec:building}
% ============================================================================

\subsection{HVAC Zone Model}

\begin{equation}
C_z\frac{\dd T_z}{\dd t} = Q_{\text{HVAC}} + Q_{\text{internal}} + Q_{\text{solar}} - UA(T_z - T_{\text{out}})
\end{equation}

Transfer function from HVAC power to zone temperature:
\begin{equation}
\frac{T_z(s)}{Q_{\text{HVAC}}(s)} = \frac{1/UA}{\tau s + 1}
\end{equation}

where $\tau = C_z/UA$.

\subsection{Lighting Control}

Dimming dynamics:
\begin{equation}
\tau_L\frac{\dd L}{\dd t} + L = K_L \cdot d
\end{equation}

where $L$ is light level, $d$ is dimmer setting, and $\tau_L$ is lamp response time.

% ============================================================================
\section{Traffic Flow}
\label{sec:traffic}
% ============================================================================

\subsection{Cell Transmission Model}

For a road segment:
\begin{equation}
\frac{\dd n_i}{\dd t} = q_{i-1} - q_i
\end{equation}

where $n_i$ is vehicle count in cell $i$ and $q_i$ is flow rate.

Flow-density relationship:
\begin{equation}
q = v_f n \left(1 - \frac{n}{n_{\max}}\right)
\end{equation}

\subsection{Traffic Signal Timing}

Green time allocation affects queue discharge:
\begin{equation}
\frac{\dd q}{\dd t} = \lambda - s \cdot g(t)
\end{equation}

where $s$ is saturation flow rate and $g(t) \in \{0, 1\}$ indicates green signal.

% ============================================================================
